% Options for packages loaded elsewhere
\PassOptionsToPackage{unicode}{hyperref}
\PassOptionsToPackage{hyphens}{url}
\PassOptionsToPackage{dvipsnames,svgnames,x11names}{xcolor}
%
\documentclass[
  10pt,
  letterpaper,
]{article}
\usepackage{amsmath,amssymb}
\usepackage{iftex}
\ifPDFTeX
  \usepackage[T1]{fontenc}
  \usepackage[utf8]{inputenc}
  \usepackage{textcomp} % provide euro and other symbols
\else % if luatex or xetex
  \usepackage{unicode-math} % this also loads fontspec
  \defaultfontfeatures{Scale=MatchLowercase}
  \defaultfontfeatures[\rmfamily]{Ligatures=TeX,Scale=1}
\fi
\usepackage{lmodern}
\ifPDFTeX\else
  % xetex/luatex font selection
\fi
% Use upquote if available, for straight quotes in verbatim environments
\IfFileExists{upquote.sty}{\usepackage{upquote}}{}
\IfFileExists{microtype.sty}{% use microtype if available
  \usepackage[]{microtype}
  \UseMicrotypeSet[protrusion]{basicmath} % disable protrusion for tt fonts
}{}
\makeatletter
\@ifundefined{KOMAClassName}{% if non-KOMA class
  \IfFileExists{parskip.sty}{%
    \usepackage{parskip}
  }{% else
    \setlength{\parindent}{0pt}
    \setlength{\parskip}{6pt plus 2pt minus 1pt}}
}{% if KOMA class
  \KOMAoptions{parskip=half}}
\makeatother
\usepackage{xcolor}
\usepackage[margin=0.78in]{geometry}
\usepackage{color}
\usepackage{fancyvrb}
\newcommand{\VerbBar}{|}
\newcommand{\VERB}{\Verb[commandchars=\\\{\}]}
\DefineVerbatimEnvironment{Highlighting}{Verbatim}{commandchars=\\\{\}}
% Add ',fontsize=\small' for more characters per line
\newenvironment{Shaded}{}{}
\newcommand{\AlertTok}[1]{\textcolor[rgb]{1.00,0.00,0.00}{\textbf{#1}}}
\newcommand{\AnnotationTok}[1]{\textcolor[rgb]{0.38,0.63,0.69}{\textbf{\textit{#1}}}}
\newcommand{\AttributeTok}[1]{\textcolor[rgb]{0.49,0.56,0.16}{#1}}
\newcommand{\BaseNTok}[1]{\textcolor[rgb]{0.25,0.63,0.44}{#1}}
\newcommand{\BuiltInTok}[1]{\textcolor[rgb]{0.00,0.50,0.00}{#1}}
\newcommand{\CharTok}[1]{\textcolor[rgb]{0.25,0.44,0.63}{#1}}
\newcommand{\CommentTok}[1]{\textcolor[rgb]{0.38,0.63,0.69}{\textit{#1}}}
\newcommand{\CommentVarTok}[1]{\textcolor[rgb]{0.38,0.63,0.69}{\textbf{\textit{#1}}}}
\newcommand{\ConstantTok}[1]{\textcolor[rgb]{0.53,0.00,0.00}{#1}}
\newcommand{\ControlFlowTok}[1]{\textcolor[rgb]{0.00,0.44,0.13}{\textbf{#1}}}
\newcommand{\DataTypeTok}[1]{\textcolor[rgb]{0.56,0.13,0.00}{#1}}
\newcommand{\DecValTok}[1]{\textcolor[rgb]{0.25,0.63,0.44}{#1}}
\newcommand{\DocumentationTok}[1]{\textcolor[rgb]{0.73,0.13,0.13}{\textit{#1}}}
\newcommand{\ErrorTok}[1]{\textcolor[rgb]{1.00,0.00,0.00}{\textbf{#1}}}
\newcommand{\ExtensionTok}[1]{#1}
\newcommand{\FloatTok}[1]{\textcolor[rgb]{0.25,0.63,0.44}{#1}}
\newcommand{\FunctionTok}[1]{\textcolor[rgb]{0.02,0.16,0.49}{#1}}
\newcommand{\ImportTok}[1]{\textcolor[rgb]{0.00,0.50,0.00}{\textbf{#1}}}
\newcommand{\InformationTok}[1]{\textcolor[rgb]{0.38,0.63,0.69}{\textbf{\textit{#1}}}}
\newcommand{\KeywordTok}[1]{\textcolor[rgb]{0.00,0.44,0.13}{\textbf{#1}}}
\newcommand{\NormalTok}[1]{#1}
\newcommand{\OperatorTok}[1]{\textcolor[rgb]{0.40,0.40,0.40}{#1}}
\newcommand{\OtherTok}[1]{\textcolor[rgb]{0.00,0.44,0.13}{#1}}
\newcommand{\PreprocessorTok}[1]{\textcolor[rgb]{0.74,0.48,0.00}{#1}}
\newcommand{\RegionMarkerTok}[1]{#1}
\newcommand{\SpecialCharTok}[1]{\textcolor[rgb]{0.25,0.44,0.63}{#1}}
\newcommand{\SpecialStringTok}[1]{\textcolor[rgb]{0.73,0.40,0.53}{#1}}
\newcommand{\StringTok}[1]{\textcolor[rgb]{0.25,0.44,0.63}{#1}}
\newcommand{\VariableTok}[1]{\textcolor[rgb]{0.10,0.09,0.49}{#1}}
\newcommand{\VerbatimStringTok}[1]{\textcolor[rgb]{0.25,0.44,0.63}{#1}}
\newcommand{\WarningTok}[1]{\textcolor[rgb]{0.38,0.63,0.69}{\textbf{\textit{#1}}}}
\usepackage{longtable,booktabs,array}
\usepackage{calc} % for calculating minipage widths
% Correct order of tables after \paragraph or \subparagraph
\usepackage{etoolbox}
\makeatletter
\patchcmd\longtable{\par}{\if@noskipsec\mbox{}\fi\par}{}{}
\makeatother
% Allow footnotes in longtable head/foot
\IfFileExists{footnotehyper.sty}{\usepackage{footnotehyper}}{\usepackage{footnote}}
\makesavenoteenv{longtable}
\setlength{\emergencystretch}{3em} % prevent overfull lines
\providecommand{\tightlist}{%
  \setlength{\itemsep}{0pt}\setlength{\parskip}{0pt}}
\setcounter{secnumdepth}{5}
\newcommand{\Z}{\mathbb Z}
\newcommand{\C}{\mathbb C}
\newcommand{\Pone}{\mathbb P^1}
\newcommand{\cO}{\mathcal O}
\newcommand{\Pic}{\operatorname{Pic}}
\newcommand{\NS}{\operatorname{NS}}
\newcommand{\rk}{\operatorname{rk}}
\newcommand{\im}{\operatorname{im}}
\newcommand{\coker}{\operatorname{coker}}
\newcommand{\Aut}{\operatorname{Aut}}
\newcommand{\Spec}{\operatorname{Spec}}
\newcommand{\dP}{\mathrm{dP}}
\ifLuaTeX
  \usepackage{selnolig}  % disable illegal ligatures
\fi
\usepackage{bookmark}
\IfFileExists{xurl.sty}{\usepackage{xurl}}{} % add URL line breaks if available
\urlstyle{same}
\hypersetup{
  pdftitle={Revised Full Mathematical Peer Review of the 108-Page S6 Manuscript},
  pdfauthor={Independent technical review},
  colorlinks=true,
  linkcolor={blue},
  filecolor={Maroon},
  citecolor={Blue},
  urlcolor={blue},
  pdfcreator={LaTeX via pandoc}}

\title{Revised Full Mathematical Peer Review of the 108-Page S6
Manuscript}
\usepackage{etoolbox}
\makeatletter
\providecommand{\subtitle}[1]{% add subtitle to \maketitle
  \apptocmd{\@title}{\par {\large #1 \par}}{}{}
}
\makeatother
\subtitle{Corrective addendum, proof-status reconstruction, and
Codex-finding integration}
\author{Independent technical review}
\date{25 August 2026}

\begin{document}
\maketitle

{
\hypersetup{linkcolor=}
\setcounter{tocdepth}{3}
\tableofcontents
}
\section{Corrective addendum dated 25 August
2026}\label{corrective-addendum-dated-25-august-2026}

This edition integrates the mathematical claims contained in the
user-supplied Codex progress log, but it does not treat that log as
evidence by itself. The log reports two strengthened proof lines: a
direct period-theoretic proof that

\[
a(X)=1,
\qquad
\mathcal M(X)=f^*\mathbb C(t),
\]

and an integral nearby-cycle/Leray computation that proves the
homology-sphere conclusion without the global cusp-collapse theorem. It
also reports corrections to a modular-frame factor, a Neron--Severi
simplification, Lemma A.3's variance, and the distinction between
\(\Lambda\) and \(\bar\Lambda\). Only the progress log, not the separate
Codex workbench files or its claimed 38-check harness, was supplied.
Every point promoted in this report was therefore re-located and
re-checked in the audited 108-page manuscript itself.

The result is a genuine but limited correction to the review of 24
August 2026.

\begin{enumerate}
\def\labelenumi{\arabic{enumi}.}
\tightlist
\item
  The direct algebraic-dimension argument in Section 9.1 is stronger
  than the earlier review stated. Its period calculation, very-general
  Neron--Severi computation, exclusion of algebraic dimensions one and
  two for the fibres, and meromorphic descent form a complete argument
  conditional only on the construction of the period family and standard
  structure theory for complex tori. It is independent of Section 7 and
  of Campana--Demailly--Peternell.
\item
  The second topology route in Section 7.7 and Appendix B is not merely
  corroborative. It is a logically sufficient topology engine for the
  recognition of \(X\) as \(S^6\), provided the based fundamental-group
  relations, the finite-fibre specialization lattices, the integral cusp
  specialization theorem, and the geometric identification of the Leray
  transgressions are correct. Proposition 7.2's global collapse is
  therefore not theorem-critical for the sphere-recognition chain,
  although it remains necessary for the manuscript's first topology
  route and for its explicit collapse-map and generator claims.
\item
  Lemma A.3 contains a confirmed variance error. The displayed
  Mayer--Vietoris sequence and every map subsequently computed are
  cohomological, although the proof writes \(H_k\). Replacing \(H_k\) by
  \(H^k\) throughout that calculation and applying universal
  coefficients at the end repairs the lemma without changing its stated
  groups.
\item
  The corrected source already contains the modular factor
  \(\mathsf e=F^{-1}e\) needed on source page 71, and the current
  algebra in Lemma 9.3 is correct. Those are repairs present in the
  audited text, not defects of the current version.
\item
  The current cusp notation correctly uses
  \(\bar\Lambda=\Lambda/\Lambda_{\mathrm{tor}}\) as the deck group. The
  Codex log does not identify the earlier erroneous occurrence, so no
  additional current-source error is asserted on that point.
\end{enumerate}

The recommendation remains \textbf{major revision and independent
specialist replication}. The reason is no longer a suspicion that the
global collapse is the sole route to the homology calculation, nor an
unresolved algebraic-dimension argument. The remaining burden is
concentrated in the actual geometric realization of the displayed local
and global maps, integral nearby-cycle functoriality and orientation
data, the based van Kampen and transgression comparisons, and the
analytic duality/base-change/literature-comparison chain in Section 10.

\subsection{Mathematical results so far: a cold-entry
guide}\label{mathematical-results-so-far-a-cold-entry-guide}

\subsubsection{The construction being
claimed}\label{the-construction-being-claimed}

The manuscript starts with a rank-four integral lattice \(\Lambda\) and
monodromies of orbifold orders \(3\) and \(4\), together with a
unipotent cusp monodromy. Period functions \(\tau,\mu,\beta\) define a
holomorphic family of two-dimensional complex tori over the
thrice-punctured projective line. Three local fillings are then
attached: a toric, reduced, irreducible, nonnormal normal-crossing fibre
\(W\) at the cusp, and multiple fibres \(3S_1\) and \(4S_2\) whose
reductions are bielliptic surfaces. The gluing parameters enter the
topology only through

\[
p=12\ell_0-4\ell_1-3\ell_2.
\]

For the actual data \((\ell_0,\ell_1,\ell_2)=(0,1,-1)\), one has
\(p=-1\).

\subsubsection{The topology chain now
isolated}\label{the-topology-chain-now-isolated}

The based van Kampen calculation produces

\[
\pi_1(X)=\langle c,x,y\mid c\ \text{central},\ xy=c^{\ell_0},\ x^3=c^{\ell_1},\ y^4=c^{\ell_2}\rangle
\cong \mathbb Z/|p|,
\]

conditional on the geometric identification of the three meridians and
their drift terms. Thus the actual value \(p=-1\) gives \(\pi_1(X)=1\).

The independent Leray route uses the exact cusp-specialization statement

\[
H^q(W;\mathbb Z)\xrightarrow{\ \operatorname{sp}^q\ }
H^q(F;\mathbb Z)^{T_0}
\]

as an isomorphism for \(0\le q\le4\), together with the finite-fibre
specialization indices. It obtains

\[
H^0(B,R^1f_*\mathbb Z)=12\mathbb Z\gamma,
\quad
H^0(B,R^2f_*\mathbb Z)=\mathbb Z\cdot2q,
\quad
H^0(B,R^3f_*\mathbb Z)=\mathbb Z\cdot2\gamma uw.
\]

The three relevant transgressions are

\[
\begin{aligned}
d_2^{0,1}(12\gamma)&=\pm p\,\omega,\\
d_2^{0,2}(2q)&=\pm p\,\omega\bar\delta,\\
d_2^{0,3}(2\gamma uw)&=\pm p\,\omega[\gamma\delta],
\end{aligned}
\]

with one common sign. Multiplicativity gives the same integer in all
three formulas. When \(p=-1\), these maps are units. The spectral
sequence then yields \(H^1(X)=H^2(X)=H^3(X)=0\); universal coefficients
and Poincare duality give the integral homology of \(S^6\). Together
with \(\pi_1(X)=1\) and \(\Theta_6=0\), this implies
\(X\cong_{\mathrm{diff}}S^6\).

This route does not use the global collapse theorem. Its
theorem-critical unresolved inputs are instead: the based meridian
relations, the integral specialization theorem at \(W\), the
finite-fibre lattices, and the geometric realization of the first
transgression and its multiplicative comparison.

\subsubsection{The direct algebraic-dimension
chain}\label{the-direct-algebraic-dimension-chain}

For an integral class

\[
E=E_{\gamma u}\gamma u+E_{\gamma w}\gamma w+E_{\gamma\delta}\gamma\delta
  +E_{uw}uw+E_{u\delta}u\delta+E_{w\delta}w\delta,
\]

Section 9.1 proves

\[
E\in\operatorname{NS}(F_b)
\quad\Longleftrightarrow\quad
P_E(z)=0,
\]

where

\[
P_E=E_{\gamma u}-E_{\gamma w}\tau-E_{\gamma\delta}\mu
 +6E_{uw}\mu+E_{u\delta}\beta+E_{w\delta}(6\mu^2-\tau\beta).
\]

The transformation laws under the two orbifold generators prove that

\[
1,\ \tau,\ \mu,\ \beta,\ 6\mu^2-\tau\beta
\]

are linearly independent. Hence the very general fibre has

\[
\operatorname{NS}(F_b)=\mathbb Z\eta,
\qquad
\eta=uw+6\gamma\delta.
\]

Moreover

\[
\eta^2=12,
\qquad
G_\eta=
\begin{pmatrix}
2\operatorname{Im}\tau&2\operatorname{Im}\mu\\
2\operatorname{Im}\mu&\frac13\operatorname{Im}\beta
\end{pmatrix},
\]

and

\[
\det G_\eta=\frac{2}{3}\operatorname{Im}\tau
\left(\operatorname{Im}\beta-\frac{6(\operatorname{Im}\mu)^2}{\operatorname{Im}\tau}\right)<0.
\]

Thus every nonzero multiple of \(\eta\) is nondegenerate and indefinite,
while its square is nonzero. A two-torus of algebraic dimension two
would have an ample positive class; one of algebraic dimension one would
have a nonzero fibre class of square zero. Neither can occur when
\(\operatorname{NS}=\mathbb Z\eta\). Therefore the very general fibre
has algebraic dimension zero.

A separate descent lemma shows that a meromorphic function on \(X\)
which is constant on uncountably many smooth fibres comes from the base.
Since very general fibres have algebraic dimension zero,

\[
\mathcal M(X)=f^*\mathcal M(\mathbb P^1)=f^*\mathbb C(t),
\qquad
a(X)=1,
\]

and \(f\) is the algebraic reduction. This proof is independent of the
topology and of the CDP dispute.

\subsubsection{What Section 10 has and has not
established}\label{what-section-10-has-and-has-not-established}

The local conductor calculation produces a nonzero ambient cotangent
section on a reduced nonnormal normal-crossing fibre even when the
torsion-free intrinsic cotangent sheaf has no section. For a line bundle
\(A\) trivial after normalization, the pair

\[
\bigl(\eta^*(df|_W)\otimes e,0\bigr)
\]

descends through the normalization/conductor sequence to

\[
s\in H^0(W,\Omega_X^1|_W\otimes A),
\]

whose intrinsic image is conductor-supported torsion. This local
mechanism is mathematically real and was independently reconstructed.

Its application to the global contradiction requires the exact analytic
chain

\[
H^2(W,E)^\vee
\simeq \operatorname{Hom}_W(E,\omega_W)
\simeq H^0(W,E^\vee\otimes\omega_W),
\qquad E=(T_X\otimes L)|_W,
\]

followed by the top-degree base-change map

\[
(R^2f_*(T_X\otimes L))_0\otimes k(0)
\longrightarrow H^2(W,(T_X\otimes L)|_W).
\]

The manuscript states the relevant Gorenstein duality and top-degree
base-change hypotheses, and no internal contradiction was found.
Nevertheless, three assertions must remain distinct:

\begin{enumerate}
\def\labelenumi{\arabic{enumi}.}
\tightlist
\item
  the local computation may expose a gap in a printed proof at a
  nonnormal fibre;
\item
  that does not by itself prove the published theorem false in all
  stated generality;
\item
  the theorem is contradicted by this example only after the
  construction and every hypothesis on \(X\) have independently been
  certified.
\end{enumerate}

\subsection{Revised status ledger}\label{revised-status-ledger}

\begin{longtable}[]{@{}
  >{\raggedright\arraybackslash}p{(\columnwidth - 4\tabcolsep) * \real{0.3333}}
  >{\raggedright\arraybackslash}p{(\columnwidth - 4\tabcolsep) * \real{0.3333}}
  >{\raggedright\arraybackslash}p{(\columnwidth - 4\tabcolsep) * \real{0.3333}}@{}}
\toprule\noalign{}
\begin{minipage}[b]{\linewidth}\raggedright
Status
\end{minipage} & \begin{minipage}[b]{\linewidth}\raggedright
Claim
\end{minipage} & \begin{minipage}[b]{\linewidth}\raggedright
Present assessment
\end{minipage} \\
\midrule\noalign{}
\endhead
\bottomrule\noalign{}
\endlastfoot
\textbf{VERIFIED exact algebra} & monodromy, Smith forms,
invariant/coinvariant lattices, period transformations, \(P_E\),
\(\eta^2\), Gram determinant, Appendix B rank matrices & Reproduced in
two local harnesses: 28 legacy checks and 14 supplemental checks, all
passing. \\
\textbf{CONDITIONALLY VALIDATED} & direct proof of \(a(X)=1\) and
\(\mathcal M(X)=f^*\mathbb C(t)\) & The complete argument is present and
survives line-by-line review, conditional on the period-family
construction and standard complex-torus structure theory. \\
\textbf{CONDITIONALLY VALIDATED} & Leray route from specialization and
transgression data to \(H_*(X)=H_*(S^6)\) & Formal sheaf,
spectral-sequence, rank, saturation, and multiplicativity logic survives
review. Geometric bridge inputs remain open to specialist
replication. \\
\textbf{OPEN} & existence and exact gluing of the global compact complex
threefold & No fatal defect found, but the infinite toric quotient,
finite logarithmic fillings, and Hausdorff/proper gluing remain
construction-level obligations. \\
\textbf{OPEN} & Section 10 as a counterexample to the printed CDP
conclusions & Local conductor mechanism is strong; global duality/base
change and exact comparison with the published proof require specialist
adjudication. \\
\textbf{REFUTED local wording} & Remark 7.28 as a fixed-triple
biconditional & False; must be existential in \(\ell_0\). \\
\textbf{REFUTED proof sentence} & Proposition 4.7(iv)'s properness
justification & Closedness needs the closed-graph/section argument. \\
\textbf{REFUTED notation/variance} & Lemma A.3's displayed
Mayer--Vietoris sequence as homology & It is cohomology. The stated
groups survive after correction and UCT. \\
\end{longtable}

\section{Referee determination}\label{referee-determination}

\textbf{Recommendation:} Major revision and resubmission after
independent specialist replication. The complete manuscript should not
be dismissed as arithmetically incoherent or as automatically refuted by
the printed CDP20 conclusion. It also should not yet be represented as
an established solution of the Hopf problem.

\textbf{Revised central conclusion:} No fatal internal contradiction has
been identified in the complete 108-page manuscript. The direct
algebraic-dimension theorem is now assessed as substantially validated,
conditional on the earlier period-family construction. The
nearby-cycle/Leray topology route is now assessed as a logically
sufficient alternative to the global collapse for the sphere-recognition
conclusion. These are meaningful upgrades from the 24 August review.

\textbf{What remains unproved by this audit:} Exact algebra is not the
main uncertainty. The remaining high-risk obligations are the geometric
identification of the based meridians and their drifts; the realization
of the stated local specialization maps by the actual cusp and
multiple-fibre models; the integral nearby-cycle edge map and its
supported Gysin comparison; the first Leray transgression as the same
integer \(p\) appearing in van Kampen; and the analytic
duality/base-change/literature comparison in Section 10.

The manuscript claims a compact complex threefold \(X\) and a
holomorphic map

\[
f:X\longrightarrow \Pone
\]

whose general fibres are complex two-tori, whose three special fibres
consist of a reduced irreducible nonnormal normal-crossing surface \(W\)
and two multiple fibres \(3S_1\) and \(4S_2\), and for the chosen
logarithmic twisting data satisfies

\[
\pi_1(X)=1,
\qquad H^*(X;\Z)\cong H^*(S^6;\Z),
\qquad X\cong_{\mathrm{diff}}S^6.
\]

It further claims

\[
a(X)=1,
\qquad K_X\cong f^*\cO_{\Pone}(-1)\otimes\cO_X(2S_2),
\]

\[
R^1f_*\cO_X\cong\cO\oplus\cO(-1),
\qquad R^2f_*\cO_X\cong\cO(-1),
\qquad c_3(X)=2.
\]

If correct, these claims give a complex structure on the standard smooth
six-sphere and contradict printed conclusions of Campana, Demailly, and
Peternell. The burden of proof is therefore exceptional.

The disposition has three levels.

\begin{enumerate}
\def\labelenumi{\arabic{enumi}.}
\tightlist
\item
  \textbf{Finite and formal mathematics:} The displayed integer
  matrices, exterior powers, Smith forms, invariant and coinvariant
  lattices, period identities, Neron--Severi kernel, intersection form,
  Appendix B rank bookkeeping, and corrected Lemma A.3 map have been
  independently reproduced.
\item
  \textbf{Conditional proof reconstruction:} The complete direct
  algebraic-dimension chain and the formal
  nearby-cycle/Leray-to-homology chain survive line-by-line
  reconstruction under explicitly isolated geometric inputs.
\item
  \textbf{Publication-level certification:} The global construction, the
  bridge from local models to the topological and sheaf-theoretic maps,
  and the Section 10 analytic/literature claims still require
  independent domain specialists. Failure to locate an error is not a
  certificate of an extraordinary theorem.
\end{enumerate}

\section{Scope, source status, and review
method}\label{scope-source-status-and-review-method}

\subsection{Source completeness and
provenance}\label{source-completeness-and-provenance}

The reviewed archive, \texttt{s6\_canonical\_108\_pages\_audited.zip},
is complete as a reconstruction package. It contains ten modular source
parts, a modular assembler, a monolithic source, a 108-page compiled
PDF, an authority reassembly, build scripts, checksum manifests,
page-by-page visual audits, and reconstruction provenance. The archive's
own verification script passes. The modular and monolithic sources
compile to byte-identical PDFs in the audited environment.

This establishes fidelity to the supplied authority PDF, not authorship
of the TeX source and not mathematical correctness. The archive
expressly states that it is an independent reconstruction rather than
the author's original TeX. For any formal submission, an
author-controlled source or an explicit author sign-off on the
reconstructed canonical PDF is necessary.

The earlier review package was based on an incomplete 155-line assembler
plus packet 06. Its conclusion that the main theorem was directly
contradicted by CDP20 was a correct description of the logical conflict
visible in that slice, but an incorrect final assessment of the complete
manuscript because Section 10 was unavailable. This report replaces that
review in full.

\subsection{Review procedures}\label{review-procedures}

The audit used four complementary procedures.

First, I read the complete proof dependency chain rather than reviewing
isolated claims. In particular, I traced the construction from the
monodromy representation through the period family, the three local
fillings, the global gluing, both integral topology computations, smooth
recognition, analytic invariants, and the literature-conflict analysis.

Second, I independently reproduced exact computations with integer and
symbolic arithmetic. The legacy harness contains twenty-eight passing
checks. A second supplemental harness prompted by the Codex log contains
fourteen passing checks, including the corrected modular-frame
automorphy, the current Lemma 9.3 simplification, the cohomological
Mayer--Vietoris matrix behind Lemma A.3, the two Appendix B lattice
maps, the rank-forcing argument in B.5--B.6, and the common Leray
transgression integer. These are forty-two check executions across two
harnesses, with some intentional overlap; they are not advertised as
forty-two logically independent theorems. The separate Codex claim of a
38/38 core harness is recorded but not verified because that harness was
not supplied.

Third, I reconstructed the dependency graph of the two topology routes.
The first uses an explicit collapse and Mayer--Vietoris. The second uses
integral nearby cycles and the Leray spectral sequence. The second route
does not depend on Proposition 7.2 for sphere recognition; its own
indispensable bridges are now listed separately. Agreement remains
significant, but the two routes are not treated as independent where
they share local gluing, meridian, or specialization data.

Fourth, I reviewed Section 10 as an argument about exact proof
mechanisms rather than as an appeal to exception. The review
distinguishes the local conductor lemma, its global application to \(W\)
through duality and base change, the precise comparison with the printed
CDP reductions, and the logically later assertion that the published
conclusion is falsified by the candidate.

\subsection{Meaning of audit labels}\label{meaning-of-audit-labels}

``Verified computationally'' means that the displayed finite calculation
has been reproduced exactly. ``No contradiction found'' means that the
proof was read and its logical transitions appear coherent, but it is
not a certification. ``Major revision request'' identifies a pivotal
argument whose present exposition is too compressed for an extraordinary
theorem, even where no specific counterexample has been found.
``Confirmed error'' is reserved for a statement or proof sentence for
which a definite defect can be exhibited.

\section{The proof architecture and its dependency
graph}\label{the-proof-architecture-and-its-dependency-graph}

The main theorem has the following dependency chain.

\begin{enumerate}
\def\labelenumi{\arabic{enumi}.}
\tightlist
\item
  Sections 2--3 construct a rank-four integral local system over the
  \((3,4,\infty)\) orbifold and a nonpolarized holomorphic family of
  complex two-tori with monodromies of orders \(3\), \(4\), and
  unipotent type at the cusp.
\item
  Section 4 constructs a toric filling at the cusp from an infinite
  \(A_2\) fan and a free, properly discontinuous \(\Z^2\) action. Its
  central fibre \(W\) is irreducible and nonnormal, with normalization
  \(\dP_6\), three double curves, and two triple points.
\item
  Section 5 constructs the order-\(3\) and order-\(4\) logarithmic
  fillings, with reduced fibres bielliptic surfaces.
\item
  Section 6 glues the smooth family and three fillings into a compact
  complex threefold \(X\to\Pone\) and records the discrete gluing
  parameter \(\ell_0\).
\item
  Section 7 computes \(\pi_1(X)\) and integral homology. For the chosen
  twists \((\ell_0,\ell_1,\ell_2)=(0,1,-1)\), the controlling integer \[
  p=12\ell_0-4\ell_1-3\ell_2
  \] is \(-1\). The manuscript gives two integral homology computations.
\item
  Section 8 uses the simply connected integral homology-sphere
  conclusion, Hurewicz--Whitehead theory, and \(\Theta_6=0\) to identify
  \(X\) diffeomorphically with \(S^6\).
\item
  Section 9 computes algebraic dimension, direct images, the canonical
  bundle, Chern-number consequences, holomorphic forms, and
  automorphisms.
\item
  Section 10 explains why the printed CDP98/CDP20 proof mechanisms are
  claimed not to cover this \(X\).
\end{enumerate}

A fatal defect in Sections 2--6 would invalidate the space. A defect in
either topology route could potentially be repaired by the other only if
their shared local specialization and gluing data remain valid. A defect
in Section 9 would not by itself destroy the complex structure on
\(S^6\), but it could remove the claimed contradiction with CDP20. A
defect in Section 10 would reinstate the literature obstruction unless
the published theorem itself were independently corrected.

\subsection{Conditional recognition theorem and obligation
split}\label{conditional-recognition-theorem-and-obligation-split}

The minimal logical route to the smooth-sphere conclusion can be stated
without the global collapse.

\textbf{Conditional recognition theorem.} Assume the following five
inputs.

\begin{itemize}
\tightlist
\item
  \textbf{(G)} Sections 2--6 produce a closed connected smooth
  six-manifold \(X\) with a proper map \(f:X\to\mathbb P^1\) and the
  stated smooth and special fibres.
\item
  \textbf{(V1)} The based van Kampen inclusions, primitive meridians,
  signs, logarithmic drifts, and the cusp null-homotopy give the
  presentation of Theorem 7.17 with \(p=12\ell_0-4\ell_1-3\ell_2\).
\item
  \textbf{(B1)} At the cusp, integral specialization is injective with
  image \(H^q(F;\mathbb Z)^{T_0}\) in every degree.
\item
  \textbf{(F1)} At the two finite fibres, the specialization lattices
  and indices are those of Proposition 7.14 and Lemma A.4.
\item
  \textbf{(L1)} The Leray transgression \(d_2^{0,1}(12\gamma)\) is
  geometrically the obstruction class \(\pm p\omega\), and the
  multiplicative identifications on the \(E_2\) page are compatible with
  the stated invariant/coinvariant generators.
\end{itemize}

Then, for the actual data \(p=-1\), the underlying smooth manifold of
\(X\) is diffeomorphic to \(S^6\).

\textbf{Proof.} Under (V1), \(\pi_1(X)\cong\mathbb Z/|p|=1\). Under
(B1), (F1), and the already verified monodromy algebra, Proposition 7.26
gives the stated \(E_2\) page. Under (L1), multiplicativity gives all
three transgressions as multiplication by the same nonzero integer
\(p\). When \(p=-1\), the relevant kernels and cokernels vanish, so
\(H^1(X)=H^2(X)=H^3(X)=0\). Universal coefficients imply
\(H_1(X)=H_2(X)=0\); Poincare duality then gives
\(H_3(X)=H_4(X)=H_5(X)=0\), while \(H_0(X)=H_6(X)=\mathbb Z\). Hence
\(X\) is a simply connected integral homology six-sphere. Section 8's
Hurewicz--Whitehead argument makes it a homotopy sphere, and
\(\Theta_6=0\) gives the diffeomorphism. \(\square\)

The obligations attached to Proposition 7.2 should therefore be split
from this minimal set:

\begin{itemize}
\tightlist
\item
  \textbf{(C1)} a finite equivariant CW description and global
  cusp-collapse retraction;
\item
  \textbf{(C2)} the integral collapse matrix, including orientations,
  torsion, and saturation;
\item
  \textbf{(C3)} equality of collapse specialization with nearby-cycle
  specialization;
\item
  \textbf{(MV1)} geometric identification of the printed Mayer--Vietoris
  maps.
\end{itemize}

These remain necessary to certify the first topology route, the explicit
collapse map, the stated generators of \(H_*(W)\), and some later
normalizations. They are not required for the conditional theorem above.

\section{Section-by-section mathematical
review}\label{section-by-section-mathematical-review}

\subsection{Lattice and monodromy data (Section
2)}\label{lattice-and-monodromy-data-section-2}

The matrices

\[
T_1=\begin{pmatrix}
1&0&-6&2\\0&-1&1&1\\0&-1&0&1\\0&0&0&1
\end{pmatrix},\qquad
T_2=\begin{pmatrix}
1&6&0&-3\\0&0&-1&1\\0&1&0&0\\0&0&0&1
\end{pmatrix}
\]

satisfy

\[
T_1^3=I,\qquad T_2^4=I,\qquad T_0=(T_1T_2)^{-1}=I+N,\qquad N^2=0.
\]

The exact orders are indeed \(3\) and \(4\). The displayed
contragredients \(A_j=(T_j^{-1})^t\) and \(M_0=(T_0^{-1})^t\) are
correct, and \(A_1A_2M_0=I\). The invariant and vanishing-cycle lattices
are also correct:

\[
\ker(M_0-I)=\im(M_0-I)=\langle\widehat w,\widehat\delta\rangle,
\]

and the induced map

\[
B_0:\overline\Lambda\to\Lambda_{\mathrm{tor}},\qquad
B_0=\begin{pmatrix}0&1\\-1&0\end{pmatrix},
\]

is unimodular. The fixed lattices

\[
\Lambda^{A_1}=\langle\varepsilon,\widehat\delta\rangle,
\qquad
\Lambda^{A_2}=\langle\varepsilon',\widehat\delta\rangle
\]

are primitive rank-two sublattices with the stated generators.

The invariant alternating form \(Q_0\) is also consistent. In the
ordered basis \((\gamma,u,w,\delta)\), its nonzero entries are
\(Q_0(\gamma,\delta)=1\) and \(Q_0(u,w)=6\). Direct multiplication gives
\(T_j^tQ_0T_j=Q_0\). The induced form

\[
b_N(x,y)=Q_0(x,Ny)
\]

on \(V/\ker N\), in the basis \(\bar w,\bar\delta\), is
\(\operatorname{diag}(6,-1)\). This confirms the manuscript's central
nonpolarized feature: the monodromy-compatible alternating form has
indefinite limiting form and does not define a polarization. This
matters because standard polarized degeneration theorems cannot simply
be imported (Clemens, 1977; Mumford, 1972).

The group-theoretic uniformization by the \((3,4,\infty)\) triangle
group is standard. The manuscript carefully separates the uniformizing
coordinate from the period coordinate at the cusp and derives the
translation sign. I found no inconsistency in the presentation
\(\Delta\cong\Z/3*\Z/4\) or the definition of the representation.

\textbf{Disposition:} Exact algebra verified. No correction required,
except that a concise table of all actions on \(V\), \(\Lambda\), and
their exterior powers would reduce later sign risk.

\subsection{Period family (Section 3)}\label{period-family-section-3}

The period matrix is built from

\[
\sigma_1=6\mu\gamma+\tau u+w,
\qquad
\sigma_2=\beta\gamma+\mu u+\delta,
\]

with \(F^1=\langle\sigma_1,\sigma_2\rangle\). The construction proceeds
through three nontrivial analytic steps: a modular function \(\tau\), a
torsor solution \(\mu\) for an automorphy sheaf of degree \(-1\), and a
solution \(\beta\) to a further additive cocycle. The final imaginary
constant is chosen so that the lattice remains nondegenerate.

The modular transformation laws used later are mutually consistent. The
automorphy sheaf calculation in Lemma 3.10, including the degree \(-1\)
conclusion, is coherent: the chosen modular expression has the required
character, and its divisor behavior at the cusp gives
\(\mathcal L\simeq\cO(-1)\). Because a square-root factor is used, the
argument would be clearer if formulated invariantly in terms of the
orbifold line bundle and its character, rather than by an implicitly
chosen branch. I did not find a contradiction in the present
construction.

The nondegeneracy argument is important. It is not enough that
\(\Pi(z)\) be formally equivariant; its four real columns must span a
full lattice in \(\C^2\). The manuscript reduces the determinant to an
inequality involving

\[
D=\operatorname{Im}\beta-\frac{6(\operatorname{Im}\mu)^2}{\operatorname{Im}\tau}
\]

and uses the free imaginary constant in \(\beta\) to force \(D<0\)
throughout a compact fundamental domain. The proper-discontinuity and
compactness input makes this strategy plausible. The same expression
later controls the signature of the Hermitian form associated with
\(\eta\).

The cusp asymptotic decomposition

\[
Z=sB_0+C(t_c)
\]

is the precise bridge to the toric filling. I checked the equivariance
and the exponential identity that produces the torus element
\(c_{\bar\lambda}(t_c)\). No algebraic inconsistency was found.

\textbf{Disposition:} No fatal defect found. For publication, the torsor
existence and the global nondegeneracy choice should be stated as
standalone propositions with explicit cohomology groups and norm
estimates. The current proof is credible but densely packed.

\subsection{Toric filling at the cusp (Section
4)}\label{toric-filling-at-the-cusp-section-4}

The fan consists of cones over the cells of the periodic \(A_2\)
triangulation at height one. Every maximal cone is unimodular because
the edge differences form a \(\Z^2\) basis. Local finiteness of the
triangulation gives a Hausdorff, second-countable toric complex
manifold, and the character \(t=\chi^{(0,0,1)}\) has reduced
normal-crossing divisor

\[
t^{-1}(0)=\bigcup_{v\in\Z^2}E_v.
\]

The action

\[
\Psi_{\bar\lambda}=\tau_{\bar\lambda}\circ\Phi_{\bar\lambda}
\]

combines integral translation of the fan with a \(t\)-dependent torus
translation. The manuscript does more than assert properness: it
introduces a rescaled position map and obtains a quantitative
displacement estimate. Since \(B_0\) is orthogonal and unimodular, large
group elements move compact sets apart uniformly near the cusp. The
compact fundamental-set argument then yields properness of the quotient
map to the disc.

The action is free. On the central fibre, the fan translation is free on
irreducible components; away from the central fibre, the action
corresponds to a free lattice translation on the torus cover. The
quotient therefore has local models

\[
t=z_0,
\qquad t=z_0z_1,
\qquad t=z_0z_1z_2,
\]

and is smooth. The invariant logarithmic/top form gives the asserted
relative canonical trivialization.

The normalization and conductor combinatorics are consistent. A
fundamental component is \(\dP_6\); opposite sides of its toric hexagon
are identified. This yields three rational double curves and two triple
points. The orbit stratification

\[
W=(\C^*)^2\sqcup 3\C^*\sqcup\{P,Q\}
\]

gives \(e(W)=2\).

I found one definite exposition defect. Proposition 4.7(iv) says the
image of the section is closed because it is ``proper over \(\Delta\).''
That is not the correct reason when the base is an open disc. The
conclusion is readily repaired: a continuous section of a Hausdorff map
has closed graph/image, or one may prove closedness directly in the
chosen quotient chart. This does not threaten the construction.

The highest-risk part of Section 4 is not a visible algebraic error but
the global topological bookkeeping for an infinite fan and quotient. The
manuscript contains the needed ideas, yet a publication version should
consolidate them into one theorem: local finiteness, second
countability, freeness, proper discontinuity, compact fundamental
domain, and compatibility of the toric deformation with the group
action.

\textbf{Disposition:} No theorem-breaking defect found. One minor proof
sentence should be corrected; global quotient topology deserves
expansion.

\subsection{Finite-order logarithmic fillings (Section
5)}\label{finite-order-logarithmic-fillings-section-5}

The local family over the linearizing disc carries an affine
transformation whose linear part is \(A_j\) and whose translation part
is \(v_j/m_j\). The admissibility conditions ensure that the affine
cyclic action is free. The fixed-point criterion is correctly reduced to
membership of the translation vector in \((A_j-I)\) of the torus
lattice, and the chosen vectors avoid that image.

The quotient produces a smooth threefold with multiple fibre \(m_jS_j\).
The reductions \(S_j\) are bielliptic surfaces. The computed first
homology and torsion are compatible with the standard Bagnera--de
Franchis classification (Barth et al., 2004). The order-\(3\) and
order-\(4\) cases also generate the precise local specialization indices
later used in Section 7.

The logarithmic regluing section is arranged so that the transformed
zero section is invariant. The conjugation identity used to compare the
punctured family with the quotient is correct. The resulting drift
controls \(\ell_j=\gamma(v_j)\).

\textbf{Disposition:} No fatal error found. The proof would be more
auditable with a single diagram showing the cover, affine action,
quotient, multiple fibre, reduction, and boundary identification. The
classification of the two bielliptic surfaces should be tied explicitly
to the standard notation in Barth et al.~(2004).

\subsection{Global gluing and compactness (Section
6)}\label{global-gluing-and-compactness-section-6}

The three local models are glued to the smooth family over \(B^\circ\).
The overlap maps are biholomorphisms by construction. The manuscript
addresses Hausdorffness through shrinkable collars and the fact that
each local piece maps properly to a disc. Properness of the global map
is local on the base for a separated holomorphic map; the resulting
total space is compact because the base is compact and the map is
proper.

The zero section is used as more than notation: it fixes the lift of
boundary meridians and prevents hidden fibre translations in the
finite-order gluings. At the cusp, affine regluing changes the discrete
parameter \(\ell_0\), and the tautological gluing gives \(\ell_0=0\).

No incompatibility was found among the boundary identifications. The
proof is nevertheless unusually sensitive to the distinction between the
smooth fibre lattice, the toric translation lattice, and the quotient of
\(\Lambda\) by \(\Lambda_{\mathrm{tor}}\). A one-page gluing table
should identify, for every boundary component, the deck group, meridian,
fibre translation, orientation, and induced relation in \(\pi_1\).

\textbf{Disposition:} Coherent. Major revision only in presentation and
auditability, not because of a demonstrated flaw.

\subsection{Fundamental group (Section
7.5)}\label{fundamental-group-section-7.5}

After attaching the cusp filling, the normal closure of
\(\Lambda_{\mathrm{tor}}\) is \(\ker\gamma\). Hence the image of the
fibre group is generated by one central element \(c\). The two finite
fillings impose

\[
x^3=c^{\ell_1},\qquad y^4=c^{\ell_2},
\]

and the cusp relation gives \(xy=c^{\ell_0}\). Thus

\[
\pi_1(X)=\langle c,x,y\mid c\text{ central},\ xy=c^{\ell_0},\ x^3=c^{\ell_1},\ y^4=c^{\ell_2}\rangle.
\]

Eliminating \(y\) makes the group abelian with relation matrix

\[
R=\begin{pmatrix}
3&-\ell_1\\
4&\ell_2-4\ell_0
\end{pmatrix},
\qquad
\det R=-(12\ell_0-4\ell_1-3\ell_2).
\]

Because the entries include \(3\) and \(4\), the first Smith invariant
is \(1\); hence

\[
\pi_1(X)\cong\Z/|p|,
\qquad
p=12\ell_0-4\ell_1-3\ell_2.
\]

The admissibility congruences imply \(\gcd(p,12)=1\). For \((0,1,-1)\),
\(p=-1\) and the group is trivial. For the comparison sign, \(p=-7\).

These calculations pass exact Smith-normal-form checks. The manuscript's
long sign lemma is not redundant: a relative sign error would
interchange the simply connected and \(\Z/7\) examples. I found its
treatment internally consistent. The final invariant \(|p|\) is
orientation-independent, but the manuscript should provide a compact
convention table because the proof currently requires the reader to
track clockwise meridians through several covers.

\textbf{Disposition:} Algebra verified. No mathematical correction
required.

\subsection{First integral homology route: collapse and Mayer--Vietoris
(Sections
7.1--7.6)}\label{first-integral-homology-route-collapse-and-mayervietoris-sections-7.17.6}

The first route constructs an equivariant collapse of a nearby torus
fibre onto the singular fibre and uses Mayer--Vietoris to obtain

\[
H_k(W;\Z)=(\Z,\Z^2,\Z^4,\Z^2,\Z),
\qquad
H_k(\mathcal M;\Z)\text{ of ranks }(1,3,6,6,3,1),
\]

followed by the cyclic groups \(\mathbb Z/|p|\) in degrees \(1\) through
\(4\) for \(X\). The exterior-power matrices, saturated image lattices,
rank of \(\alpha_2\), and primitive left kernel all pass exact checking.

The Codex-directed reinspection changes two aspects of this section.

First, Proposition 7.2 is no longer classified as a bottleneck for
sphere recognition because Section 7.7 and Appendix B bypass it. It
remains a high-risk theorem for the manuscript's explicit collapse
description and first homology route. The appropriate obligations are
(C1)--(C3) and (MV1), not an undifferentiated request to ``check the
topology.''

Second, Lemma A.3 has a confirmed variance error. The displayed sequence

\[
\cdots\to H_k(W)\to H_k(\widetilde W)\oplus H_k(D)
\to H_k(\widetilde D)\to H_{k+1}(W)\to\cdots
\]

has the direction of the \textbf{cohomological} Mayer--Vietoris
sequence. The subsequent map is restriction

\[
H^2(\widetilde W)\oplus H^2(D)\longrightarrow H^2(\widetilde D),
\]

computed by intersection numbers on \(\mathrm{dP}_6\) and pullback from
the three double curves. Its integral matrix has rank \(4\) and Smith
diagonal \((1,1,1,1)\); hence kernel rank \(3\) and free cokernel rank
\(2\). The preceding connecting homomorphism
\(H^1(\widetilde D)\to H^2(W)\) contributes one additional primitive
\(\mathbb Z\). Thus

\[
H^\bullet(W;\mathbb Z)=(\mathbb Z,\mathbb Z^2,\mathbb Z^4,\mathbb Z^2,\mathbb Z),
\]

all free, and universal coefficients recovers the same homology groups.
The lemma's conclusion survives; its written proof must be corrected.

\textbf{Disposition:} Exact finite algebra passes. The route is
conditionally coherent after the Lemma A.3 variance repair. Its global
collapse and map-identification claims remain specialist-verification
items, but they are not part of the minimal recognition theorem.

\subsection{Second integral homology route: Leray spectral sequence
(Section
7.7)}\label{second-integral-homology-route-leray-spectral-sequence-section-7.7}

This route is now assessed as an independent minimal topology engine
rather than merely as corroboration. It depends on the actual geometric
inputs (B1), (F1), and (L1), but it does not depend on the global
collapse of Section 7.2.

Let \(j:B^\circ\hookrightarrow B\) and let \(\mathcal L^q\) have stalk
\(\bigwedge^qV\). The cusp theorem identifies \(R^qf_*\mathbb Z\) with
\(j_*\mathcal L^q\) near \(p_0\). The finite fibres contribute the
quotient lattices

\[
Q^1=\mathbb Z/3\oplus\mathbb Z/4,
\qquad
Q^2=Q^3=\mathbb Z/2\ \text{at }p_2,
\qquad
Q^4=\mathbb Z/3\oplus\mathbb Z/4.
\]

The resulting global sections and coinvariants are

\[
\begin{aligned}
H^0(B,R^1)&=12\mathbb Z\gamma,& H^2(B,R^1)&=\mathbb Z[\bar\delta],\\
H^0(B,R^2)&=\mathbb Z\cdot2q,& H^2(B,R^2)&=\mathbb Z[\gamma\delta],\\
H^0(B,R^3)&=\mathbb Z\cdot2\gamma uw,& H^2(B,R^3)&=\mathbb Z[uw\delta].
\end{aligned}
\]

The parabolic groups \(H^1(B,R^1)\) and \(H^1(B,R^2)\) vanish by
explicit lattice-intersection calculations. These intersections and
saturation statements have been reproduced exactly.

The first transgression is not merely a symbolic parameter. It is
computed by choosing lifts over the three special discs and over a
pair-of-pants neighbourhood, evaluating their discrepancies on clockwise
meridians, and using

\[
\widetilde a_0=t_\mu\hat\rho_2^{-1}\hat\rho_1^{-1}.
\]

The sum of the three discrepancy values is \(-p\), giving

\[
d_2^{0,1}(12\gamma)=\pm p\omega.
\]

The remaining transgressions are determined by the Leibniz rule. If

\[
d_2(2q)=p'\omega\bar\delta,
\qquad
d_2(2\gamma uw)=p''\omega[\gamma\delta],
\]

then the relation \((12\gamma)(2q)=12(2\gamma uw)\) gives

\[
p''=2p-p'.
\]

The vanishing of \((2q)(2\gamma uw)\) and the two primitive pairings
give \(|p'|=|p''|\). Since \(p\ne0\) under admissibility,

\[
(p')^2-(p'')^2=4p(p'-p)=0
\]

forces \(p'=p''=p\). This exact argument passes the supplemental
harness.

For \(p=-1\) the three maps are units. The spectral-sequence deductions,
universal-coefficient step, and Poincare-duality step are correct. The
unresolved question is therefore sharply geometric: do the actual local
models realize the specialization lattices and based discrepancy class
exactly as asserted? This is the combined content of (V1), (B1), (F1),
and (L1).

\textbf{Disposition:} Conditionally validated as a complete topology
route. Independent specialist replication should target the geometric
bridges, not repeat the already verified Smith and spectral-sequence
arithmetic.

\subsection{Recognition as the smooth six-sphere (Section
8)}\label{recognition-as-the-smooth-six-sphere-section-8}

Once \(X\) is simply connected and has the integral homology of \(S^6\),
the Hurewicz theorem and Whitehead's theorem make it a homotopy
\(6\)-sphere. Smooth homotopy \(6\)-spheres are classified by
\(\Theta_6\), and \(\Theta_6=0\) (Kervaire \& Milnor, 1963). Thus the
diffeomorphism conclusion follows.

This step is standard and correct, conditional on Section 7. The
manuscript could make explicit that the smooth structure is oriented by
the complex orientation, although in dimension six there is no exotic
sphere ambiguity.

\textbf{Disposition:} Valid conditional deduction.

\subsection{Algebraic dimension (Section
9.1)}\label{algebraic-dimension-section-9.1}

This section is upgraded from ``no contradiction found'' to
\textbf{conditionally validated}. The proof is complete in the audited
manuscript, and it is independent of Section 7 and of CDP20.

For \(E\in H^2(F_b;\mathbb Z)\), Lemma 9.2 proves that \(E\) is of type
\((1,1)\) precisely when

\[
E\wedge\sigma_1\wedge\sigma_2=0,
\]

which is the equation \(P_E(z)=0\) with

\[
P_E=E_{\gamma u}-E_{\gamma w}\tau-E_{\gamma\delta}\mu
+6E_{uw}\mu+E_{u\delta}\beta+E_{w\delta}(6\mu^2-\tau\beta).
\]

The wedge expansion and all signs have been reproduced. Lemma 9.3 uses
the \(g_0\) and \(g_2\) transformation laws to show that

\[
1,\tau,\mu,\beta,6\mu^2-\tau\beta
\]

are linearly independent. The two subtraction identities, including the
corrected algebraic simplification reported by Codex, are correct in the
current source. Consequently

\[
P_E\equiv0
\quad\Longleftrightarrow\quad
E\in\mathbb Z\eta,
\qquad
\eta=uw+6\gamma\delta.
\]

For every smooth fibre, \(\eta\) is a \((1,1)\) class with
\(\eta^2=12\). The Hermitian form has determinant

\[
\frac23\operatorname{Im}\tau
\left(\operatorname{Im}\beta-\frac{6(\operatorname{Im}\mu)^2}{\operatorname{Im}\tau}\right)<0,
\]

so every nonzero multiple is nondegenerate and indefinite. Outside a
countable set, \(\operatorname{NS}(F_b)=\mathbb Z\eta\). The
complex-torus algebraic-reduction theorem then gives a direct dichotomy:

\begin{itemize}
\tightlist
\item
  algebraic dimension two would supply an ample class with
  positive-definite Hermitian form;
\item
  algebraic dimension one would supply the nonzero pullback of a point
  from an elliptic quotient, whose square is zero.
\end{itemize}

Both contradict the properties of \(n\eta\). Therefore the very general
fibre has algebraic dimension zero.

Lemma 9.6 supplies the global descent. If a meromorphic function is
constant on uncountably many fibres, the analytic set

\[
Z=\{dg\wedge f^*dt=0\}
\]

must be all of the complement of its polar divisor over \(B^\circ\);
otherwise its countably many irreducible components could contain only
countably many distinct fibre components. The function is therefore
fibrewise constant and descends holomorphically away from finitely many
base points. A transverse analytic disc gives polynomial growth and
hence meromorphic extension across each missing point. Thus

\[
\mathcal M(X)=f^*\mathbb C(t),
\qquad a(X)=1,
\]

and \(f\) is the algebraic reduction.

The remaining conditions are upstream rather than internal to this
proof: the period functions must define the claimed family and satisfy
the nondegeneracy inequalities, and the standard structure theorem for
algebraic reduction of complex tori must be cited in its exact form. No
Section 7 topology and no CDP theorem is used.

\textbf{Disposition:} Conditionally validated; revise the manuscript
mainly by making the standard torus theorem and the
connectivity/irreducibility facts in the descent lemma
citation-explicit.

\subsection{Canonical bundle and direct images (Sections
9.2--9.4)}\label{canonical-bundle-and-direct-images-sections-9.29.4}

The relative two-form \(e=d\zeta_1\wedge d\zeta_2\) transforms by
\(\det R_g\). Its behavior at the finite orbifold points gives
fractional vanishing orders \(2/3\) and \(1/4\), and it is regular and
nonvanishing at the cusp. The resulting direct image is

\[
f_*\omega_{X/B}\cong\cO_{\Pone}(1).
\]

The multiple-fibre correction yields

\[
K_X\cong f^*\cO(-1)\otimes\cO(2S_2).
\]

The character arithmetic is coherent. The non-torsion conclusion follows
because the stated bundle has a nontrivial multiple-fibre component and
its appropriate power is a nontrivial pullback.

For the structure-sheaf direct images, the manuscript obtains

\[
f_*\cO_X=\cO,
\qquad
R^1f_*\cO_X\cong\cO\oplus\cO(-1),
\qquad
R^2f_*\cO_X\cong\cO(-1).
\]

The fibre cohomology at \(W\) is computed by the normalization/conductor
sequence. The finite multiple fibres are handled by equivariant local
models and descending induction through the nilpotent thickening.
Relative duality identifies the top direct image. The eigenweights
determine the splitting of \(R^1\).

I found no algebraic inconsistency, and the resulting global Hodge
numbers agree with the Leray calculation. The main concern is the
precision of the analytic foundations. In a theorem of this importance,
references to Grauert, base change, and relative duality should be
attached to exact statements and hypotheses. The reader should not have
to infer that the scheme-theoretic multiple fibres are covered, that the
relevant sheaves are flat, and that duality applies to the Gorenstein
nonnormal surface \(W\).

\textbf{Disposition:} Plausible and internally consistent. Major
revision requested for exact analytic theorem statements.

\subsection{Holomorphic forms, automorphisms, and Chern data (Sections
9.5--9.7)}\label{holomorphic-forms-automorphisms-and-chern-data-sections-9.59.7}

The vanishing of global holomorphic \(p\)-forms is derived by writing
invariant forms on the universal cover and using the period derivatives
and local growth at the cusp and orbifold points. The weight-zero and
weight-one lemmas eliminate the remaining coefficients. The argument is
technically detailed and consistent with the computed \(h^{0,q}\).

The identity component of the automorphism group must preserve the
algebraic reduction and hence acts fibrewise. Translation vector fields
correspond to the invariant lattice line
\(\Lambda^\Delta=\Z\widehat\delta\), giving a one-dimensional vertical
field. Its flow is \(\C^*\) rather than an elliptic translation group
because of the degeneration. The toric fixed set is one rational double
curve with normal weights \((+1,-1)\). This is compatible with
\(e(X^{\C^*})=e(X)=2\).

Because \(H^2(X;\Z)=H^4(X;\Z)=0\), \(c_1\) and \(c_2\) vanish
integrally. The top Chern number equals the Euler characteristic,

\[
\langle c_3(T_X),[X]\rangle=e(X)=2.
\]

The Riemann--Roch arithmetic giving \(\chi(\cO_X)=0\),
\(\chi(\Omega_X^1)=-1\), and \(\chi(T_X)=1\) is consistent.

\textbf{Disposition:} No contradiction found. These conclusions remain
conditional on the topology and direct-image calculations.

\section{Appendix B: integral specialization at the nonnormal
cusp}\label{appendix-b-integral-specialization-at-the-nonnormal-cusp}

Appendix B is the central bridge in the collapse-free topology proof.
The Codex log correctly identifies it as the place where exact matrices
alone are insufficient. A line-by-line reconstruction nevertheless
supports a stronger disposition than the previous report: the formal
nearby-cycle, hypercohomology, rank, saturation, and supported-Gysin
chain is coherent and complete once the stated local toric models and
standard integral nearby-cycle functoriality are granted.

\subsection{Nearby-cycle sheaves and their integral
gluing}\label{nearby-cycle-sheaves-and-their-integral-gluing}

The local equations \(t=z_0\), \(z_0z_1\), and \(z_0z_1z_2\) give Milnor
fibres of homotopy type point, \(S^1\), and \(T^2\). Hence

\[
\mathcal H^0\psi=\mathbb Z_W,
\qquad
\mathcal H^1\psi=\mathcal V,
\qquad
\mathcal H^2\psi=\mathbb Z_P\oplus\mathbb Z_Q.
\]

The normalization globally distinguishes the two branches along each
double curve. With the manuscript's orientation convention, the three
restrictions at both triple points satisfy

\[
n_1-n_2+n_3=0.
\]

This gives

\[
0\to\mathcal V\to\bigoplus_{k=1}^3\iota_{k*}\mathbb Z_{\bar C_k}
\xrightarrow{\delta_1}\mathbb Z_P\oplus\mathbb Z_Q\to0,
\qquad
\delta_1=
\begin{pmatrix}1&-1&1\\1&-1&1\end{pmatrix}.
\]

The matrix has rank one and primitive image. The claimed cohomology

\[
H^\bullet(W;\mathcal V)=(\mathbb Z^2,\mathbb Z,\mathbb Z^3,0)
\]

follows integrally.

\subsection{The bottom row and the Lemma A.3
repair}\label{the-bottom-row-and-the-lemma-a.3-repair}

The normalization/conductor resolution of \(\mathbb Z_W\) yields the map

\[
d:H^2(\mathrm{dP}_6;\mathbb Z)\to\mathbb Z^3,
\qquad
d(aH-\textstyle\sum b_iE_i)=(s,-s,s),
\quad s=b_1+b_2+b_3-a.
\]

Its image is primitive. Thus

\[
H^\bullet(W;\mathbb Z)=(\mathbb Z,\mathbb Z^2,\mathbb Z^4,\mathbb Z^2,\mathbb Z),
\]

all free. This is the cohomological calculation that Lemma A.3 writes
with the wrong variance. Appendix B itself uses the correct cohomology
direction.

\subsection{Rank forcing and the degree-three supported
class}\label{rank-forcing-and-the-degree-three-supported-class}

The \(E_2\) ranks are

\[
\begin{array}{c|ccccc}
b=2&2&0&0&0&0\\
b=1&2&1&3&0&0\\
b=0&1&2&4&2&1\\ \hline
&a=0&a=1&a=2&a=3&a=4.
\end{array}
\]

Writing the five possible differential ranks as
\((\alpha,\beta,\gamma,\delta,\epsilon)\), total-degree ranks force

\[
\alpha=\gamma=0,
\qquad
\beta+\delta+\epsilon=1.
\]

The supported argument takes a compact \(T_f\)-orbit \(K\cong T^2\) in
\(W^\circ\), transports it smoothly to a translate \(K_t\) of the
saturated linear subtorus in \(F_t\), and compares

\[
H^3(W,W\setminus K)\to H^3(W)\xrightarrow{\operatorname{sp}^3}H^3(F_t)
\]

with the Gysin map

\[
H^1(K_t)\xrightarrow{\iota_!}H^3(F_t).
\]

Because \(F_t\cong T_f\times T_B\) topologically,
\(\iota_!(a)=a\times[T_B]^\vee\) is injective with saturated image of
rank two. Hence \(\operatorname{sp}^3\) has rank two, forcing

\[
(\alpha,\beta,\gamma,\delta,\epsilon)=(0,0,0,1,0).
\]

The remaining filtration quotients are free, so the specialization
images are saturated. Their ranks equal the ranks of the
\(T_0\)-invariants, and inclusion into the invariant lattice therefore
becomes equality.

The supplemental harness reproduces the two integral matrices and the
complete rank-forcing logic.

\subsection{Exact residual
obligations}\label{exact-residual-obligations}

The formal proof is sound under the following items, which should be
citation- and convention-complete in a revision.

\begin{enumerate}
\def\labelenumi{\arabic{enumi}.}
\tightlist
\item
  The local quotient charts of Theorem 4.5 must really be
  holomorphically equivalent to the three stated normal-crossing
  equations, with the stated branch labels and compact-torus action.
\item
  The specialization morphism \(\mathbb Z_W\to\psi_f\mathbb Z\) and its
  hypercohomology edge map must agree with the sheaf specialization used
  in Section 7.7 over the integers.
\item
  The smooth supported transport in Lemma B.3 must be tied functorially
  to that edge map. The map \(\Theta\) is a smooth open embedding, which
  is sufficient for singular cohomology; it should not be read as
  holomorphic.
\item
  The marking must identify the transported \(T_f\)-orbit with the
  saturated sublattice \(\Lambda_{\mathrm{tor}}\) and hence with the
  primitive Kunneth summand used by the Gysin map.
\item
  The first Leray transgression must be compared with the based
  meridian/gluing data without hidden sign or translation changes.
\end{enumerate}

These are precise bridge lemmas, not missing matrix computations.

\textbf{Disposition:} Conditionally validated as an integral
specialization proof, with specialist replication required for the five
geometric/functorial inputs above.

\section{Section 10 and the conflict with
CDP98/CDP20}\label{section-10-and-the-conflict-with-cdp98cdp20}

\subsection{Four distinct propositions that must not be
conflated}\label{four-distinct-propositions-that-must-not-be-conflated}

The literature issue decomposes into four propositions.

\begin{enumerate}
\def\labelenumi{\arabic{enumi}.}
\tightlist
\item
  \textbf{Local proposition:} a reduced nonnormal normal-crossing fibre
  can carry a conductor-supported ambient cotangent section invisible in
  the torsion-free intrinsic cotangent sheaf.
\item
  \textbf{Example-specific analytic proposition:} the line bundles
  arising from \(X\) satisfy the normalization, duality, and top-degree
  base-change hypotheses needed to turn that section into
  \(R^2f_*(T_X\otimes L)\ne0\).
\item
  \textbf{Proof-comparison proposition:} the printed CDP reduction uses
  an implication invalidated by this local phenomenon, or otherwise
  fails to account for it.
\item
  \textbf{Counterexample proposition:} the fully constructed \(X\)
  satisfies every hypothesis of the printed theorem and violates its
  conclusion.
\end{enumerate}

The first proposition is independently reconstructed and strong. The
second appears plausible and internally coherent but requires exact
analytic references. The third requires specialist comparison with the
original proof. The fourth additionally depends on certification of the
entire construction and topology. A possible defect in a proof, falsity
of a theorem, and correctness of the candidate are logically separate.

\subsection{The formal contradiction}\label{the-formal-contradiction}

If the manuscript's main conclusions hold, then \(X\) satisfies

\[
H^1(X;\Z)=H^2(X;\Z)=0,
\qquad a(X)=1,
\]

and has a holomorphic algebraic reduction \(f:X\to\Pone\). Campana et
al.~(2020, Theorem 2.2) print, under those hypotheses, \(c_3(X)\le0\),
and their Corollary 2.3 prints \(a(X)=0\) for a compact complex
threefold homeomorphic to \(S^6\). The manuscript instead has
\(c_3(X)=2\) and \(a(X)=1\). Both sets of conclusions cannot
simultaneously be correct.

The complete manuscript does not evade this issue. It isolates two exact
proof mechanisms and argues that one is irreparably false at the
nonnormal fibre \(W\), while the other is a monodromy error that can be
repaired for this example.

\subsection{The nonnormal-fibre cotangent
mechanism}\label{the-nonnormal-fibre-cotangent-mechanism}

Let \(S\) be a reduced normal-crossing fibre in a smooth threefold
\(\mathcal X\to\Delta\), with normalization \(\eta:\widetilde S\to S\)
and conductor \(D\). Locally the defining function is \(g=ux\), \(uxy\),
or \(uxyz\). The differential \(dg|_S\) is nonzero away from the
conductor but restricts to zero on the conductor.

For a line bundle \(A\) on \(S\) whose pullback to \(\widetilde S\) is
trivial, choose a nowhere-zero pullback section \(e\). The
normalization/conductor Mayer--Vietoris sequence lets the pair

\[
\bigl(\eta^*(dg|_S)\otimes e,0\bigr)
\]

descend to a nonzero section

\[
s\in H^0(S,\Omega^1_{\mathcal X}|_S\otimes A).
\]

Its image in \(\Omega_S^1\otimes A\) is supported on the conductor and
is torsion. Therefore its image in the torsion-free quotient

\[
\widetilde\Omega_S^1=\Omega_S^1/\mathcal T(\Omega_S^1)
\]

is zero. If \(A\) is nontrivial, \(s\) does not arise from the conormal
line. This produces exactly the type of ambient cotangent section that
is invisible to a vanishing statement for
\(\widetilde\Omega_S^1\otimes A\).

This is the strongest part of the manuscript's rebuttal. The local
formulas are elementary and the descent through the conductor sequence
is explicit. The mechanism shows that the implication

\[
H^0(S,\widetilde\Omega_S^1\otimes A)=0
\quad\Longrightarrow\quad
H^0(S,\Omega^1_X|_S\otimes A)=0
\]

is not valid for nonnormal \(S\) merely from vanishing of the conormal
contribution. For normal \(S\), extension across codimension two can
restore the argument; for the nonnormal surface \(W\), the conductor
supports the obstruction.

For the constructed \(W\), every line bundle from \(X\) pulls back
trivially to the normalization \(\dP_6\) because \(c_1=0\) and
\(\Pic(\dP_6)\hookrightarrow H^2(\dP_6;\Z)\). Serre--Grothendieck
duality then gives

\[
H^2\bigl(W,(T_X\otimes L)|_W\bigr)^*
\cong
H^0\bigl(W,\Omega_X^1|_W\otimes(L^*\otimes\omega_X)|_W\bigr)
\ne0.
\]

Top-degree cohomology and base change imply
\(R^2f_*(T_X\otimes L)\ne0\). The manuscript therefore claims that the
hypothesis used by Campana et al.~(2020, Proposition 2.4) holds for no
\(L\) on this \(X\).

I found no local algebra error in this mechanism. The exact analytic
statement required is the Serre--Grothendieck duality isomorphism for
the Cartier, hence Gorenstein, compact surface \(W\):

\[
H^2(W,E)^\vee\simeq\operatorname{Hom}_W(E,\omega_W)
\simeq H^0(W,E^\vee\otimes\omega_W),
\]

with \(E=(T_X\otimes L)|_W\) and \(\omega_W\simeq\omega_X|_W\) by
adjunction for the principal reduced fibre. The required base-change
statement is the top-degree map

\[
(R^2f_*(T_X\otimes L))\otimes k(c)
\longrightarrow H^2(f^*(c),(T_X\otimes L)|_{f^*(c)}),
\]

for every scheme-theoretic fibre, including the multiple fibres. The
top-degree character makes the claim plausible, but the revision should
cite a theorem whose hypotheses explicitly cover the proper flat
analytic map and locally free sheaf used here. Because the application
would invalidate a published theorem's conclusion for this example, this
chain must be independently refereed by experts in nonnormal complex
spaces, analytic duality, and cohomology and base change.

\subsection{The monodromy-coinvariant
correction}\label{the-monodromy-coinvariant-correction}

The second issue concerns a split group extension

\[
1\to K\to G\to Q\to1.
\]

The correct abelianization is

\[
G^{\mathrm{ab}}\cong Q^{\mathrm{ab}}\oplus(K^{\mathrm{ab}})_Q,
\]

where the second summand consists of coinvariants under conjugation. It
is not generally \(Q^{\mathrm{ab}}\oplus K^{\mathrm{ab}}\) unless the
action is trivial. This is standard group homology.

For the manuscript's torus fibration, the unipotent monodromy acts
nontrivially on \(H_1(F;\Z)\). The coinvariant rank is one, not four.
Consequently, formulas in the printed CDP98/CDP20 route that use
\(b_1(F)\) rather than the coinvariant rank do not apply to this
example. This also explains how

\[
H^1(X,\cO_X)\cong\C
\]

can coexist with

\[
H^1(F,\cO_F)\cong\C^2
\]

without surjectivity of restriction.

The manuscript then proves by a separate countability argument that a
line bundle non-torsion on every fibre nevertheless exists. Thus the
monodromy problem breaks the printed route to that input but can be
repaired for this example. By contrast, the nonnormal cotangent
reduction remains unavailable.

The group-theoretic correction is correct. Its application to the exact
sequence in CDP98 should still be checked against the precise local
topological model and notation in the original paper, but the
mathematical principle is not in doubt.

\subsection{Referee assessment of the literature
dispute}\label{referee-assessment-of-the-literature-dispute}

The earlier packet-only report treated CDP20 as an external theorem that
immediately refuted the manuscript. That treatment is no longer
defensible after reading Section 10. The manuscript supplies a concrete
countermechanism, not a generic assertion of exception.

The appropriate present conclusion is narrower:

\begin{itemize}
\tightlist
\item
  The manuscript and the printed CDP20 conclusions are formally
  incompatible.
\item
  Section 10 gives a coherent explanation of why the printed proof does
  not cover the proposed nonnormal fibre.
\item
  This review found no internal defect in that explanation.
\item
  The explanation has not, by virtue of this review, been externally
  adjudicated. Acceptance would effectively publish a counterexample to
  CDP20 and should involve independent referees familiar with the
  original proof.
\end{itemize}

\section{Confirmed errors, repairs, and reported prior
corrections}\label{confirmed-errors-repairs-and-reported-prior-corrections}

\subsection{Remark 7.28: quantifier
error}\label{remark-7.28-quantifier-error}

The statement

\[
|p|=1
\quad\Longleftrightarrow\quad
(\ell_1\bmod3,\ell_2\bmod4)\in\{(1,3),(2,1)\}
\]

is false as a biconditional for a fixed triple
\((\ell_0,\ell_1,\ell_2)\). For example,

\[
(\ell_0,\ell_1,\ell_2)=(100,1,-1)
\]

has the listed residue pair but \(p=1199\). The correct statement is
existential in \(\ell_0\): for fixed admissible \((\ell_1,\ell_2)\),
there exists an integer \(\ell_0\) with \(|p|=1\) exactly for those two
residue pairs.

\subsection{Proposition 4.7(iv): closed-image
justification}\label{proposition-4.7iv-closed-image-justification}

The section image is not shown closed by calling the section ``proper
over the open disc.'' A section of a continuous map into a Hausdorff
space has closed graph, and here \(f_0\circ\varsigma=\operatorname{id}\)
gives the required closed-embedding argument directly. The conclusion is
correct; the stated justification should be replaced.

\subsection{Lemma A.3: variance error}\label{lemma-a.3-variance-error}

The displayed Mayer--Vietoris sequence and the maps used in the proof
are cohomological. Replace every \(H_k\) in that calculation by \(H^k\),
retain the rank and saturation computation, and then invoke universal
coefficients to obtain the asserted homology. The supplemental exact
matrix check confirms rank \(4\), primitive image, kernel rank \(3\),
and free cokernel rank \(2\).

\subsection{Corrections reported by Codex and present in the audited
source}\label{corrections-reported-by-codex-and-present-in-the-audited-source}

The current source uses the invariant frame \(\mathsf e=F^{-1}e\) on
source page 71. Since both \(F\) and \(e\) transform by
\(\jmath_g^{-1}\), the inverse factor is necessary and the current
expression is invariant. The current version of Lemma 9.3 also has the
correct second-difference identity. These are not current-source errors.

The cusp deck group is written as
\(\bar\Lambda=\Lambda/\Lambda_{\mathrm{tor}}\), of rank two, while
\(\Lambda\) remains the rank-four homology lattice of a smooth fibre.
That typing is consistent in the audited passages inspected. Because the
Codex log does not supply the location of an earlier
\(\Lambda/\bar\Lambda\) mistake, the review records the repair without
inventing an additional confirmed defect.

\section{Major revision requirements}\label{major-revision-requirements}

\subsection{1. Publish a typed bridge-lemma
ledger}\label{publish-a-typed-bridge-lemma-ledger}

Every central claim should be tagged as exact algebra, standard theorem,
geometric identification, analytic comparison, or literature comparison.
The minimum topology set is (G), (V1), (B1), (F1), and (L1). Each should
have a standalone statement, hypotheses, proof location, orientation
conventions, and a falsification test.

\subsection{2. Separate the collapse route from the minimal recognition
route}\label{separate-the-collapse-route-from-the-minimal-recognition-route}

Proposition 7.2 and the first Mayer--Vietoris route should be developed
under (C1)--(C3) and (MV1). The paper should state prominently that the
sphere-recognition theorem can instead use Section 7.7 and Appendix B.
Conversely, claims about the explicit collapse map, named generators of
\(H_2(W)\), and equality with nearby-cycle specialization must not be
inferred from the Leray route without proof.

\subsection{3. Make the nearby-cycle and transgression bridges
citation-complete}\label{make-the-nearby-cycle-and-transgression-bridges-citation-complete}

State the precise integral nearby-cycle functor, the specialization
morphism, the edge-map convention, the supported transport map, and the
Gysin orientation. Then give a commuting diagram comparing the first
Leray transgression with the based meridian obstruction. Exact matrix
agreement is not a substitute for this diagram.

\subsection{4. Make analytic duality and top-degree base change
theorem-complete}\label{make-analytic-duality-and-top-degree-base-change-theorem-complete}

For Proposition 9.13 and Theorem 10.5, state the morphism, sheaf,
flatness, fibre dimension, Gorenstein or Cohen--Macaulay hypothesis, and
the exact base-change map. Treat the reduced nonnormal fibre and the
nonreduced multiple fibres separately. A stable citation must cover the
analytic category actually used.

\subsection{5. Present the CDP comparison as an independently refereed
note}\label{present-the-cdp-comparison-as-an-independently-refereed-note}

Separate the local conductor proposition, the example-specific
duality/base-change application, the monodromy-coinvariant correction,
and the comparison with exact numbered statements in CDP98/CDP20. A
proof gap should not be rhetorically converted into theorem falsity
before the candidate is certified.

\subsection{6. Add a complete convention and gluing
table}\label{add-a-complete-convention-and-gluing-table}

List the lattice, quotient lattice, monodromy action, deck group, based
meridian, orientation, affine translation, logarithmic drift, and
induced relation at each special point. This table is the principal
defense against a hidden sign or lattice-typing error.

\subsection{7. Correct all three confirmed local
defects}\label{correct-all-three-confirmed-local-defects}

The revision ledger must explicitly correct Proposition 4.7(iv), Remark
7.28, and Lemma A.3. It should also identify the version in which the
page-71 modular factor, Lemma 9.3 simplification, and
\(\Lambda/\bar\Lambda\) notation were repaired.

\subsection{8. Strengthen source and proof
provenance}\label{strengthen-source-and-proof-provenance}

The reviewed archive is a reconstruction. A formal submission needs
author-controlled source or author sign-off on the canonical PDF.
Distinguish ``candidate proof,'' ``independent repair or alternative,''
and ``jointly surviving version'' at the theorem level. The supplied
Codex log is useful provenance, but unsupplied files and harnesses
cannot be treated as independently audited evidence.

\section{Computational audit}\label{computational-audit}

Two executable harnesses accompany this edition.

The legacy harness performs twenty-eight exact checks, all passing. It
covers monodromy orders and contragredients, invariant and coinvariant
lattices, Smith normal forms, Wang ranks, saturation of exterior-power
images, the fundamental-group relation matrix, period transformations,
the Neron--Severi polynomial, \(\eta^2\), and the Hermitian Gram
determinant.

The supplemental Codex-integration harness performs fourteen checks, all
passing:

\begin{itemize}
\tightlist
\item
  invariance of the corrected page-71 frame \(F^{-1}e\);
\item
  the current \(g_0\) difference and corrected second-difference algebra
  in Lemma 9.3;
\item
  the exact kernel \(P_E\equiv0\Longleftrightarrow E\in\mathbb Z\eta\);
\item
  non-isotropy of every nonzero multiple of \(\eta\);
\item
  the corrected cohomological Mayer--Vietoris matrix in Lemma A.3, with
  Smith diagonal \((1,1,1,1)\);
\item
  the primitive normalization-difference map in Appendix B;
\item
  the primitive \(\delta_1\) map defining \(\mathcal V\);
\item
  all rank candidates in B.5 and the fact that the supported rank-two
  class forces \(\delta=1\);
\item
  the factorization that forces \(p'=p''=p\) in Proposition 7.27;
\item
  unit transgressions for the actual value \(p=-1\);
\item
  the fixed-triple counterexample to Remark 7.28;
\item
  the rank consistency of \(\bar\Lambda=\Lambda/\Lambda_{\mathrm{tor}}\)
  as the cusp deck group.
\end{itemize}

There are therefore forty-two passing check executions across the two
harnesses. Several deliberately overlap the same formulas under
different logical decompositions. The number should not be read as a
measure of theorem probability.

The Codex log separately reports a 38/38 ``independent core harness.''
That code and its output were not supplied, so this review neither
adopts nor contradicts the number. It records only independently
rerunnable artifacts in the present ZIP.

No harness proves that a displayed matrix is the actual map induced by
the geometric construction. Every computational result in this package
is explicitly conditional on that identification.

\section{Final recommendation}\label{final-recommendation}

I recommend \textbf{major revision and resubmission}, followed by
independent specialist replication.

The evidential balance is now more favorable than in either the
packet-only review or the first full-manuscript review. Two major
subarguments have survived a more exact reconstruction:

\begin{enumerate}
\def\labelenumi{\arabic{enumi}.}
\tightlist
\item
  Section 9.1 gives a direct and logically independent proof of
  \(a(X)=1\) and \(\mathcal M(X)=f^*\mathbb C(t)\).
\item
  Section 7.7 with Appendix B gives a collapse-free route from the local
  specialization and transgression data to the integral homology sphere.
\end{enumerate}

The manuscript therefore should not be rejected on the grounds that
algebraic reduction remains open, that Proposition 7.2 is the only
topology engine, or that CDP20 can be cited without engaging Section 10.

Acceptance would still be premature. The candidate is extraordinary, and
the unresolved tasks lie precisely where formal algebra offers the least
protection: construction of the global analytic space, based gluing and
meridian data, integral nearby-cycle functoriality in the actual toric
quotient, geometric identification of Leray transgressions, and the
analytic and bibliographic adjudication of the CDP conflict. The three
confirmed local errors are repairable and do not presently destroy the
main chain, but they show why a typed morphism and variance ledger is
necessary.

The appropriate publication path is coordinated review by specialists
in: (1) toric or semi-abelian degenerations and infinite locally finite
fans; (2) integral nearby cycles, supported specialization, and topology
of nonnormal normal-crossing degenerations; and (3) compact complex
threefolds, analytic duality/base change, and the CDP98/CDP20 arguments.
At least one referee should reconstruct the minimal recognition theorem
from (G), (V1), (B1), (F1), and (L1) without relying on the manuscript's
prose summary.

The present verdict is therefore neither certification nor dismissal. It
is that the manuscript now has a coherent, nontrivial, and partially
independently validated proof architecture worthy of specialist
replication, with a sharply delimited list of remaining obligations.

\section{References}\label{references}

Agricola, I., Bazzoni, G., Goertsches, O., Konstantis, P., \& Rollenske,
S. (2018). On the history of the Hopf problem. \emph{Differential
Geometry and its Applications, 57}, 1--9.

Ash, A., Mumford, D., Rapoport, M., \& Tai, Y.-S. (2010). \emph{Smooth
compactifications of locally symmetric varieties} (2nd ed.). Cambridge
University Press.

Banica, C., \& Stanasila, O. (1976). \emph{Algebraic methods in the
global theory of complex spaces}. Wiley.

Barth, W., Hulek, K., Peters, C., \& Van de Ven, A. (2004).
\emph{Compact complex surfaces} (2nd ed.). Springer.

Birkenhake, C., \& Lange, H. (2004). \emph{Complex abelian varieties}
(2nd ed.). Springer.

Bredon, G. E. (1997). \emph{Sheaf theory} (2nd ed.). Springer.

Campana, F., Demailly, J.-P., \& Peternell, T. (1998). The algebraic
dimension of compact complex threefolds with vanishing second Betti
number. \emph{Compositio Mathematica, 112}, 77--91.

Campana, F., Demailly, J.-P., \& Peternell, T. (2020). Corrigendum: The
algebraic dimension of compact complex threefolds with vanishing second
Betti number. \emph{Compositio Mathematica, 156}, 679--696.

Clemens, C. H. (1977). Degeneration of Kahler manifolds. \emph{Duke
Mathematical Journal, 44}, 215--290.

Deligne, P. (1973). Le formalisme des cycles evanescents. In P. Deligne
\& N. Katz (Eds.), \emph{Groupes de monodromie en geometrie algebrique
(SGA 7 II)} (Lecture Notes in Mathematics 340, pp.~82--115). Springer.

Dimca, A. (2004). \emph{Sheaves in topology}. Springer.

Fulton, W. (1993). \emph{Introduction to toric varieties}. Princeton
University Press.

Hatcher, A. (2002). \emph{Algebraic topology}. Cambridge University
Press.

Kashiwara, M., \& Schapira, P. (1990). \emph{Sheaves on manifolds}.
Springer.

Kervaire, M. A., \& Milnor, J. W. (1963). Groups of homotopy spheres I.
\emph{Annals of Mathematics, 77}, 504--537.

Lange, H., \& Birkenhake, C. (1999). \emph{Complex tori}. Birkhauser.

Mumford, D. (1972). An analytic construction of degenerating abelian
varieties over complete rings. \emph{Compositio Mathematica, 24},
239--272.

Ramis, J.-P., \& Ruget, G. (1970). Complexe dualisant et theoremes de
dualite en geometrie analytique complexe. \emph{Publications
Mathematiques de l'IHES, 38}, 77--91.

Ramis, J.-P., Ruget, G., \& Verdier, J.-L. (1971). Dualite relative en
geometrie analytique complexe. \emph{Inventiones Mathematicae, 13},
261--283.

Smale, S. (1961). Generalized Poincare's conjecture in dimensions
greater than four. \emph{Annals of Mathematics, 74}, 391--406.

\section{Appendix A: audit disposition by central
claim}\label{appendix-a-audit-disposition-by-central-claim}

\begin{longtable}[]{@{}
  >{\raggedright\arraybackslash}p{(\columnwidth - 6\tabcolsep) * \real{0.2500}}
  >{\raggedright\arraybackslash}p{(\columnwidth - 6\tabcolsep) * \real{0.2500}}
  >{\raggedright\arraybackslash}p{(\columnwidth - 6\tabcolsep) * \real{0.2500}}
  >{\raggedright\arraybackslash}p{(\columnwidth - 6\tabcolsep) * \real{0.2500}}@{}}
\toprule\noalign{}
\begin{minipage}[b]{\linewidth}\raggedright
Central claim
\end{minipage} & \begin{minipage}[b]{\linewidth}\raggedright
Audit disposition
\end{minipage} & \begin{minipage}[b]{\linewidth}\raggedright
Principal dependency
\end{minipage} & \begin{minipage}[b]{\linewidth}\raggedright
Required follow-up
\end{minipage} \\
\midrule\noalign{}
\endhead
\bottomrule\noalign{}
\endlastfoot
Smooth compact complex threefold \(X\) exists & No fatal defect found;
not certified & Sections 2--6, obligation (G) & Independent review of
period family, infinite toric quotient, finite fillings, and gluing \\
\(\pi_1(X)=1\) & Exact group algebra verified, conditional on based
attaching relations & (V1), Theorem 7.17 & Based meridian, primitive
generator, sign, drift, and cusp null-homotopy audit \\
\(H^*(X;\mathbb Z)=H^*(S^6;\mathbb Z)\) & Leray deduction conditionally
validated & (B1), (F1), (L1) & Replicate integral specialization and
transgression comparison \\
Explicit collapse and first MV route & Algebra verified; global geometry
open & (C1)--(C3), (MV1) & Cell-level equivariant construction and map
ledger; repair Lemma A.3 variance \\
\(X\cong_{\mathrm{diff}}S^6\) & Standard conditional consequence & (G),
(V1), (B1), (F1), (L1), \(\Theta_6=0\) & None beyond topology inputs \\
\(a(X)=1\) and \(f\) algebraic reduction & Conditionally validated by
direct period/descent proof & Sections 2--3 and 9.1 & Exact citation for
complex-torus algebraic reduction; upstream period-family validation \\
Direct images and \(K_X\) & Plausible, not fully certified & Sections
9.2--9.4 & Exact analytic base-change, relative duality, and
multiple-fibre hypotheses \\
\(c_3(X)=2\) & Conditional on homology-sphere topology & Sections 7--8 &
No separate arithmetic issue \\
Local nonnormal cotangent mechanism & Independently reconstructed &
Lemmas 10.2--10.3 & Present as standalone proposition \\
CDP20 conclusion fails for \(X\) & Substantive but externally
unadjudicated & Entire main theorem plus Theorem 10.5 and literature
comparison & Specialist review of duality, base change, and exact
published proof \\
\end{longtable}

\section{Appendix B: supersession and version
history}\label{appendix-b-supersession-and-version-history}

Three review stages should remain visible.

\begin{enumerate}
\def\labelenumi{\arabic{enumi}.}
\tightlist
\item
  The packet-06 report correctly identified the formal incompatibility
  with the printed CDP20 conclusion but lacked Section 10 and the full
  construction. Its rejection recommendation is withdrawn.
\item
  The full-manuscript report of 24 August 2026 withdrew that rejection
  and recommended major revision. Its principal overstatement was to
  leave the global collapse among the central unresolved risks for
  sphere recognition and to describe Section 9.1 more cautiously than
  its complete proof warrants.
\item
  This 25 August edition supersedes the 24 August report on those two
  points, adds the Lemma A.3 variance defect, and integrates the
  supplied Codex findings under an explicit evidence/provenance
  crosswalk. The recommendation of major revision remains unchanged.
\end{enumerate}

The 24 August PDF and Markdown source are retained under
\texttt{history/} rather than silently overwritten.

\section{Appendix C: reproducibility}\label{appendix-c-reproducibility}

The revised review bundle contains:

\begin{itemize}
\tightlist
\item
  \texttt{report/revised\_full\_mathematical\_peer\_review.pdf}, the
  primary report;
\item
  the corresponding Markdown and generated LaTeX sources;
\item
  \texttt{history/full\_mathematical\_peer\_review\_2026-08-24.pdf} and
  \texttt{.md};
\item
  \texttt{audit/issue\_ledger.csv} and \texttt{.json};
\item
  \texttt{audit/theorem\_dependency\_matrix.csv};
\item
  \texttt{audit/codex\_findings\_crosswalk.csv} and \texttt{.json};
\item
  the legacy \texttt{audit/exact\_checks.py} harness and results;
\item
  \texttt{audit/codex\_integration\_checks.py} and its JSON/log outputs;
\item
  the user-supplied Codex progress log under \texttt{evidence/codex/};
\item
  reconstruction/source-verification evidence;
\item
  \texttt{MANIFEST.sha256}, archive test output, and a review receipt.
\end{itemize}

To rerun both arithmetic audits from the bundle root:

\begin{Shaded}
\begin{Highlighting}[]
\ExtensionTok{python}\NormalTok{ audit/exact\_checks.py}
\ExtensionTok{python}\NormalTok{ audit/codex\_integration\_checks.py}
\end{Highlighting}
\end{Shaded}

The report is built from Markdown with Pandoc and pdfLaTeX, rendered at
200 dpi, visually inspected, and preflighted before packaging.

\end{document}
